\documentclass[addpoints]{exam}
% \documentclass[addpoints,12pt]{exam}

\usepackage[slovene]{babel}
\usepackage{amsfonts}
\usepackage{amssymb}
\usepackage{amsmath}
\usepackage{float}
\usepackage{graphicx}
\usepackage{enumitem}
\usepackage{color}
\usepackage[gen]{eurosym}
\usepackage{newunicodechar}
\newunicodechar{€}{\euro}


%Komentar naslednje vrstice glede na verzijo testa -- rešitve ali prostor za odgovore:
% \printanswers

% \linespread{2}


\pointsinrightmargin
\marginbonuspointname{}


\renewcommand{\solutiontitle}{\noindent\textbf{Rešitve in točkovanje:}\par\noindent}

\colorgrids
\definecolor{GridColor}{gray}{0.7}
\setlength{\gridsize}{7mm}

\extrawidth{5mm}
\setlength{\rightpointsmargin}{1.5cm}

\begin{document}

\runningfooter{}{}{}

\begin{center}
    \fbox{\fbox{\parbox{15cm}{\centering
    Peto pisno ocenjevanje znanja -- ponovno \\ 1. letnik -- test B \\ 4. junij 2026 \\ (Koordinatni sistem v ravnini, Funkcija, Premica)}}}
\end{center}

\vspace{0.1in}
\makebox[\textwidth]{Ime in priimek, oddelek:\enspace\hrulefill}


\begin{center}
    \chqword{Naloga}
    \chpword{Možne točke}
    \chbpword{Dodatne točke}
    \chsword{Dosežene točke}
    \chtword{Skupaj}  
    % \combinedpointtable[h][questions]
    \combinedgradetable[h][questions]

\end{center}

\vspace{0.1in}


\begin{questions}

    % 1. vprašanje

    \question
    Narišite dano množico točk v ravnini.
    \begin{parts}
        \part[4]
        $\{-1,2,3\} \times [-1,2)$

            \begin{solution}[7cm]
                ????
            \end{solution}

        \part[4]
        $\left\{(x,y)\in\mathbb{R}^2 \mid (x<2)\land (y=x+1)\right\} $

            \begin{solution}[5.5cm]
                ????
            \end{solution}

    \end{parts}
    
    
    \newpage

    % 2. vprašanje
    \question[5]
    Določite vrednost parametra $b$ tako, da bo imel trikotnik $\triangle ABC$, 
    z oglišči v točkah $A(-3,4)$, $B(b,5)$ in $C(1,-4)$, ploščino $30$ in negativno orientacijo.
    
    
    \begin{solution}[6cm]
        ???
    \end{solution}

        


    % 3. vprašanje
    \question
    Dana je premica z enačbo $y=\frac{3}{2}x-3$, ki seka koordinatni osi v točkah $M$ in $N$.

    \begin{parts}
        \part[4]
        Narišite premico.

            \begin{solution}[9cm]
                ???
            \end{solution}

        % \part[4]
        % Določite implicitno obliko enačbe pravokotnice k dani premici, ki poteka skozi točko $(4,3)$. % Premico tudi narišite.

        %     \begin{solution}[5cm]
        %         ???
        %     \end{solution}

        \part[5]
        Naj bo $S$ razpolovišče daljice $MN$. Izračunajte razdaljo med točko $S$ in koordinatnim izhodiščem.

            \begin{solution}[2cm]
                ???
            \end{solution}

    \end{parts}
    
    \newpage


    % 5. vprašanje
    \question
    Dana je funkcija $$f(x)=\begin{cases}
    1.5x-3, & x\leq 2 \\ 2, & 2<x<5
    \end{cases}.$$ 
    
    \begin{parts}
        \part[5]
        Narišite dano funkcijo.

            \begin{solution}[11cm]
                ???
            \end{solution}

        \part[1]
        Določite zalogo vrednosti.
        
            \begin{solution}[1cm]
                ???
            \end{solution}

        \part[2]
        Izračunajte ničle in začetno vrednost funkcije.

            \begin{solution}[6cm]
                ???
            \end{solution}

        % \part[2]
        % Določite in utemeljite, ali je funkcija surjektivna. 
        
        %     \begin{solution}[3cm]
        %         ???
        %     \end{solution}

        \part[1]
        Določite interval, kjer je funkcija strogo naraščajoča.
        
            \begin{solution}[0cm]
                ???
            \end{solution}

    \end{parts}
    
    


    \newpage 


    % 4. vprašanje
    \question
        Določite vrednost parametra $c$, da bo premica $4x-\frac{1}{c}y+2=0$:
        \begin{parts}
            \part[2] 
            potekala skozi točko $(2,4)$,

                \begin{solution}[7cm]
                    ???
                \end{solution}


            \part[3]
            padajoča,
            
                \begin{solution}[5cm]
                    ???
                \end{solution}

            % \part[3] 
            % imela diferenčni količnik enak začetni vrednosti,
                        
            %     \begin{solution}[7cm]
            %         ???
            %     \end{solution}

                
            \bonuspart[3]
            \textit{Dodatna naloga:} na ordinatni osi določala odsek velikosti $6$.

        \end{parts}
    





\end{questions}



\end{document}